\documentclass[a4paper,11pt]{article}
\usepackage[utf8]{inputenc}
\usepackage[dvipdfmx]{graphicx}
\usepackage{amsmath, amssymb}

\title{空間結合QLDPC符号 \\ (Spatially-Coupled QLDPC Codes)}
\author{Siyi Yang and Robert Calderbank}
\date{}

\begin{document}

\maketitle

\begin{abstract}
空間結合（SC）符号は、その高い性能と低遅延デコーダとの親和性から、古典符号理論において十分に研究されてきた畳み込みLDPC符号のクラスである。我々は、トーリック符号を古典的な2次元空間結合（2D-SC）符号の量子的対応物として記述し、その一般化として空間結合量子LDPC（SC-QLDPC）符号を導入する。畳み込み構造を用いて、2D-SC符号のパリティ検査行列を2つの不定元を持つ多項式として表現し、2D-SC符号がスタビライザー符号となるための必要十分な代数的条件を導出する。この代数的枠組みは、新しい符号族の構築を容易にする。本論文の主眼ではないが、メモリサイズが小さいことは量子ビットの物理的な接続を容易にし、局所的なエンコードや低遅延のウィンドウデコードを可能にすることに留意されたい。本論文では、この代数的枠組みを用いて、いずれかの要素符号の短いサイクルに起因する2D-SCハイパーグラフ積（HGP）符号のタナーグラフ内の短いサイクルを最適化する。先行研究がレート1/10未満のQLDPC符号に焦点を当てているのに対し、我々は小さなメモリ、高いレート（約1/3）、および優れた閾値を持つ2D-SC HGP符号を構築する。
\end{abstract}

\section{はじめに (Introduction)}
量子コンピュータは、古典的コンピュータよりも特定の問題を解決する能力が高いことが期待されており、フォールトトレラント量子計算の理論は、エラー率が任意に低い量子コンピュータを、欠陥のある部品からどのように構築できるかを説明している。量子エラー訂正は不可欠な構成要素であり、スタビライザー符号のクラスは広範に研究されてきた。スタビライザー符号はパウリ群の可換部分群の固定部分空間であり、CSS（Calderbank-Shor-Steane）符号は、生成元を厳密にX型とZ型の演算子に分離できる特定のタイプのスタビライザー符号である。

スタビライザーの幾何学的局所性は超伝導量子ビットプラットフォームにおいて重要であり、この設計原則はトポロジカル符号の開発を促進してきた。Kitaevは、2次元スピン格子上で定義されるトポロジカル符号のクラスであるトーリック符号を導入した。Bravyiらはこれらの符号を修正し、ユークリッド平面における量子ビットの物理的な配置を可能にした。表面符号は、閉曲面上のタイリングによって定義されるトポロジカル符号の広範なクラスであり、曲面の穴が論理演算子に対応し、種数が符号のレートを決定する。DelfosseとBreuckmannは、トーリック符号のオーバーヘッドと論理エラー率を改善するために、閉じた双曲面のタイリングを使用する双曲型および半双曲型表面符号を導入した。

表面符号は非常に低いレートしか持たないが、量子LDPC（QLDPC）符号はスタビライザーの局所性を保ちつつ高いレートの実現を可能にする。TillichとZemorは、表面符号のタナーグラフが2つの反復符号のグラフ積であることを認識した上で、ハイパーグラフ積符号（2つの古典的LDPC符号のハイパーグラフ）を導入した。ハイパーグラフ積符号の最小距離はブロック長の平方根に比例して大きくなる。Leverrierは、ハイパーグラフ積符号の3次元への拡張としてXYZ符号を開発した。ハイパーグラフ積符号は、要素符号のより一般的な「積」から構築された符号の特別なケースと見なすことができる。

\subsection*{主な貢献 (Main Contributions)}
我々の主な貢献を以下に強調する：

\begin{itemize}
    \item \textbf{空間結合QLDPC符号:} トーリック符号を古典的な2次元空間結合（2D-SC）符号の量子対応物として記述し、一般化として空間結合QLDPC符号を導入する。SC-QLDPC符号のパリティ検査行列は成分行列に組み立てられ、それらが垂直に連結されてレプリカとなり、これらのレプリカが畳み込みフレームワークにおいて結合される。符号のメモリはレプリカ内の成分行列の数によって決定され、小さなメモリは量子ビットの物理的接続を単純化し、局所的なエンコードや低遅延のウィンドウデコードを可能にする。
    
    \item \textbf{特性関数 (Characteristic Functions):} 2次元畳み込みの言語を用いて、2D-SC符号のパリティ検査行列を2つの不定元 $U, V$ の多項式 $F(U,V)$ として表現する（係数はパウリ行列の行列）。我々は、2D-SC符号がスタビライザー符号であるための必要十分条件となる単純な代数的条件を導出し、この代数的枠組みによって新しい符号族の構築を容易にする。
    
    \item \textbf{最適化フレームワーク (Optimization Framework):} タナーグラフ内の短いサイクルは、ウォーターフォール領域とエラーフロア領域の両方においてBPデコーダに問題を引き起こす。擬アベル群上のリフト積符号が有限次元のSC-HGP符号と見なせるという事実は、それらが古典領域で開発された手法を用いて短いサイクルの数に関して最適化できることを意味している。一例として、いずれかの要素符号の短いサイクルから生じる2D-SC-HGP符号のタナーグラフ内の短いサイクルを最小化するための最適化フレームワークを提供する。
    
    \item \textbf{高レート量子LDPC符号 (High Rate Quantum LDPC Codes):} これまでの研究の関心は、最小距離とブロック長の線形スケーリングを達成する符号の構築や、表面符号などの低レート符号の性能評価に集中していた。我々が構築するSC-QLDPC符号は、高いレート（$1/10$未満ではなく$1/3$程度）と高い性能を兼ね備えている。
\end{itemize}

\section{準備 (Preliminaries)}
本節では、古典的な空間結合（SC）LDPC符号の構築方法およびスタビライザー符号の定義について簡潔に導入する。以降、$\mathbb{N}, \mathbb{N}^*$ はそれぞれ非負整数および正整数を表すものとする。$\mathbb{F}_2$ は2元のガロア体とする。$d$ を法とする剰余を $(\cdot)_d$ で表す。$m \times n$ の全零行列および全一行列をそれぞれ $0_{m \times n}, 1_{m \times n}$ と表記する。行列 $A$ の $(i,j)$ 成分を $(A)_{i,j}$ で表す。行列のクロネッカー積を $\otimes$ で表す。

\subsection{古典SC-LDPC符号}
バイナリ行列 $H^P$ が与えられたとき、準巡回（QC-）LDPC符号のパリティ検査行列は、$H^P$ の各要素 $(H^P)_{i,j}$ を $z \times z$ のブロックで置き換えることで構築される。各ブロックは、$(H^P)_{i,j}=0$ のときは零行列であり、そうでないときは式(1)で定義される巡回行列のべき乗となる。$H^P$ に関連付けられたタナーグラフはQC符号のプロトグラフと呼ばれ、QC符号はプロトグラフからリフトされた（持ち上げられた）と言われる。
\begin{equation}
\sigma_{z}=\begin{bmatrix}0 & \cdots & 0 & 1 \\ 1 & \cdots & 0 & 0 \\ \vdots & \ddots & \vdots & \vdots \\ 0 & \cdots & 1 & 0\end{bmatrix}.
\end{equation}

QC-SC符号は、3つの $\gamma \times \kappa$ 行列 $\Pi, P, L$ によって決定される。分割行列 $P$ は、ベース行列 $B$ をプロトグラフの $m+1$ 個の成分行列 $H_i^P$ にどのように分割するかを指定する。リフティング行列 $L$ は、QC-SC符号において $H_i^P$ がどのように成分行列 $H_i$ へとリフトされるかを指定する。ここで $i=0, \dots, m$ である。

成分行列は垂直に連結されてレプリカとなり、これらのレプリカが水平に結合される。SC符号は、パリティ検査行列の形態に応じて、テールバイティング（TB）または非テールバイティング（NTB）に分類される。SC符号のメモリは $m$ であり、レプリカの数は結合長と呼ばれ $L$ で表される。

\subsection{2次元空間結合符号}
メモリ $m_1$ のSC符号において、各成分行列 $H_i$（$i=0, \dots, m_1$）自体がメモリ $m_2$ のSC符号（成分行列 $H_{i,j}, j=0, \dots, m_2$ を持つ）である場合、その符号を2次元（2D）SC符号と呼ぶ。

\subsection{パウリ群}
$2 \times 2$ のエルミート行列は、以下の4つの単一量子ビットパウリ行列の実行列結合として一意に表現できる。
\begin{equation}
I=\begin{bmatrix}1 & 0 \\ 0 & 1\end{bmatrix}, \quad
X=\begin{bmatrix}0 & 1 \\ 1 & 0\end{bmatrix}, \quad
Y=\begin{bmatrix}0 & -i \\ i & 0\end{bmatrix}, \quad
Z=\begin{bmatrix}1 & 0 \\ 0 & -1\end{bmatrix}.
\end{equation}
$N$ 量子ビットのパウリ群 $\mathcal{P}_N$ は、これらのテンソル積に位相因子 $i^k$ を乗じたものの集合である。スタビライザー符号において重要なのは、2つのパウリ演算子が交換するかどうかを判定するシンプレクティック内積である。

\subsection{スタビライザー符号}
スタビライザー群 $S$ を、パウリ群 $\mathcal{P}_N$ の可換な部分群であって、各要素がエルミートであり、$-I$ を含まないものとして定義する。$S$ が $r$ 個の独立な生成元を持つとき、対応するスタビライザー符号は、$S$ のすべての要素によって不変な量子状態の空間として定義される。これは $K=N-r$ 個の論理量子ビットを $N$ 個の物理量子ビットに符号化するため、$[[N, K]]$ スタビライザー符号と呼ばれる。


\section{SC符号としてのトーリック符号 (Toric Codes as SC codes)}

トーリック符号は、$d \times d$ の格子上に定義されるCSS符号である。図1において、エッジは量子ビットを表し、面に接続するエッジはX型の生成元を、頂点に接続するエッジはZ型の生成元を表す。格子の幾何学的構造により、任意のスタビライザー生成元のペアが可換であることが導かれる。

我々は、$3 \times 3$ のトーリック符号のパリティ検査行列 $H$ の行を整理し、$H$ の構造が式(20)に示すように2D-SC符号の構造と類似するようにした。図1のラベルに基づくと、列 $i$ はエッジ $i$ に関連付けられた量子ビットを表す。行 $2i-1$ は面 $i$ に関連付けられたスタビライザー生成元を表し、行 $2i$ は頂点 $i$ に関連付けられた生成元を表す。例えば、15行目は $X_{10} \otimes X_{13} \otimes X_{14} \otimes X_{15}$ を表し、10行目は $Z_{2} \otimes Z_{3} \otimes Z_{4} \otimes Z_{7}$ と面8を表す。

\begin{equation}
A = \begin{bmatrix} X & I \\ I & I \end{bmatrix}, \quad
B = \begin{bmatrix} X & X \\ Z & I \end{bmatrix}, \quad
C = \begin{bmatrix} I & X \\ Z & Z \end{bmatrix}, \quad
D = \begin{bmatrix} I & I \\ I & Z \end{bmatrix}.
\end{equation}

\noindent \textbf{注釈1（$d \times d$格子への拡張）}: レベル $i$ における面/頂点のラベルは、レベル $i-1$ における面/頂点のラベルに $d$ を加え、右に巡回させることで得られる。したがって、面とその面の左上にある頂点は同じラベル1を受け取る（図1と同様である）。正確に言えば、$(i, j)$ 番目の面（$i, j = 1, \dots, d$）は $((i-1)d+(j-i)_d+1)$ とラベル付けされる。ここで $(\cdot)_d$ は $d$ を法とする剰余を表す。$(i, j)$ 番目の面の左側と下側のエッジは、それぞれ連続して $2((i-1)d+(j-i-1)_d)+1$ および $2((i-1)d+(j-i-1)_d)+2$ とラベル付けされる。$d \times d$ のトーリック符号のパリティ検査行列 $H$ は式(9)の形式をとり、ここで $A, B, C, D$ は式(19)で与えられたブロック行列である。これは、$(m_1, m_2, L_1, L_2) = (1, 1, d, d)$ のパラメータを持つ2D-SC符号の構造を有している。

\section{可換代数 (Commutative Algebra)}

我々はまず、すべての成分が零行列 $0_2$ または $2 \times 2$ のパウリ行列であるような行列に対して、シンプレクティック内積を拡張することから始める。$m_1 \times n$ 行列 $P=[P_{i,k}]$ と $m_2 \times n$ 行列 $Q=[Q_{j,k}]$ が与えられたとき、$\langle P,Q \rangle_s$ を以下の式で与えられる $m_1 \times m_2$ バイナリ行列として定義する：
\begin{equation}
(\langle P,Q \rangle_s)_{i,j} = \langle \bigoplus_{k: P_{i,k} \neq 0_2} P_{i,k}, \bigoplus_{k: Q_{j,k} \neq 0_2} Q_{j,k} \rangle_s = \sum_{\substack{k: P_{i,k} \neq 0_2 \\ \text{and } Q_{j,k} \neq 0_2}} \langle P_{i,k}, Q_{j,k} \rangle_s. \label{eq:symplectic_matrix}
\end{equation}

\textbf{定理1.} $\mathcal{P}$ を単一量子ビット上のパウリ群とする。\\
パラメータ $(m_1, m_2, L_1, L_2)$ を持つ2D-SC符号の成分行列を
\[
  H_{i,j} \in \mathcal{P}^{r \times n} \quad (i=0,\dots,m_1, \ j=0,\dots,m_2)
\]
とする。以下の条件が満たされる場合、その2D-SC符号はスタビライザー符号である：
\begin{align}
\sum_{i=0}^{m_1-d_1} \sum_{j=0}^{m_2-d_2} \langle H_{i,j}, H_{i+d_1, j+d_2} \rangle_s &= 0_{r \times r}, \quad \forall 0 \le d_1 \le m_1, \ 0 \le d_2 \le m_2, \label{eq:thm1_cond1} \\
\sum_{i=0}^{m_1-d_1} \sum_{j=0}^{m_2-d_2} \langle H_{i+d_1, j}, H_{i, j+d_2} \rangle_s &= 0_{r \times r}, \quad \forall 0 \le d_1 \le m_1, \ 0 \le d_2 \le m_2. \label{eq:thm1_cond2}
\end{align}

\noindent \textbf{注釈4.} 古典の世界では、NTB（非テールバイティング）符号は結合鎖の両端に低次数の検査ノードが存在するため、TB（テールバイティング）符号よりも優れた閾値を持つことが知られている。しかし、量子の世界では、スタビライザー符号に要求される可換性条件のため、NTB符号は最小距離が低くなるという問題に直面する。

\subsection{量子ハードウェア上のSC符号}
量子ビットの接続性に対する制限は、様々な量子ハードウェアプラットフォーム上でQLDPC符号を実装する際の大きな課題である。しかし、SC符号の畳み込み構造は、多様な量子アーキテクチャ上でそれらを効率的に採用することを可能にする。
超伝導プラットフォーム（IBMやGoogleなど）は、隣接する量子ビット間のエンタングルメントのみを許容するため、任意の接続を持つQLDPC符号を実装するには大量のSWAPゲートが必要となる。短いメモリを持つSC符号は、これらのプラットフォームにおいて、平面的なサブグラフへの効率的な分割（厚さを抑えること）を可能にする。
また、QuEraやPASQALなどの中性原子配列（再構成可能な原子配列：RAA）においても、SC符号は異なるレプリカのアンシラ量子ビットを同一レイアウトのAOD配列に分散させることで、高い並列性を自然にサポートし、実装の忠実度を向上させる。

\subsection{特性関数 (Characteristic Function)}
我々は2次元畳み込みの言語を用いて、定理1の条件を構築と解析を容易にするシンプルな代数形式へと変換する。\\
\textbf{定義2.} 2D-SC符号の特性関数は、以下の多項式で定義される：
\begin{equation}
F(U,V) = \sum_{i=0}^{m_1} \sum_{j=0}^{m_2} H_{i,j} U^i V^j. \label{eq:characteristic_func}
\end{equation}

\textbf{定理2.} 2D-SC符号がスタビライザー符号であるための必要十分条件は、特性関数 $F(U,V)$ が以下を満たすことである：
\begin{equation}
\langle F(U,V), F(U^{-1}, V^{-1}) \rangle_s = 0_{r \times r}. \label{eq:thm2_cond}
\end{equation}

\noindent \textbf{例6（2D-SC HGP符号）.}
行列多項式 $A(U,V) \in \mathbb{F}_2^{r_1 \times n_1}[U,V]$ および $B(U,V) \in \mathbb{F}_2^{r_2 \times n_2}[U,V]$ が与えられたとき、以下を定義する：
\begin{equation}
F(U,V) = \begin{bmatrix} X(I_{n_2 \times n_2} \otimes A(U,V)) & X(\overline{B}(U,V)^T \otimes I_{r_1 \times r_1}) \\ Z(B(U,V) \otimes I_{n_1 \times n_1}) & Z(I_{r_2 \times r_2} \otimes \overline{A}(U,V)^T) \end{bmatrix}.
\end{equation}
我々は、$\langle F(U,V), F(U^{-1}, V^{-1}) \rangle_s = 0$ が成り立つことを確認することで、この2D-SCハイパーグラフ積符号がスタビライザー符号であることを示せる。

\noindent \textbf{注釈6（擬アベル・リフト積符号は有限次元SC-HGP符号である）.} 
上記の構成が、リフト積符号の表現と類似していることに注目されたい。すべての有限生成アベル群が巡回群の直積であることを利用すると、擬アベル・リフト積符号（LP符号）は有限次元SC-HGP符号と見なすことができる。これは、古典的なSC-LDPC符号のために開発されたサイクル最適化手法を通じて、高性能な擬アベルLP符号を効率的に構築できることを意味している。

\section{符号の最適化 (Code Optimization)}
本章では、4.2節で導入されたSC-HGP符号を最適化するためのフレームワークを開発する。ベース行列あるいはプロトグラフ内の閉路（closed path）が、いつプロトグラフまたはタナーグラフ内の閉路やサイクルを生み出すかを規定する「交替和条件（alternating sum condition）」から始める。その後、短いサイクルの数を効率的に削減する最適化アルゴリズムを開発する。

\subsection{サイクル条件 (Cycle Conditions)}
古典符号理論において、タナーグラフ内の短いサイクルは、ウォーターフォール領域およびエラーフロア領域の両方において、確率伝搬（BP）デコーダに問題を引き起こす。タナーグラフ内のサイクルの数を減らすことはLDPC符号の最適化における主要な目標であるが、一般的なLDPC符号においてサイクルの数を効率的に減らすことは容易な作業ではない。Progressive edge growth (PEG) は、タナーグラフの周囲長（girth：サイクルの最小長）を効率的に改善するアルゴリズムである。

しかし、サイクルの多重度も重要であるため、周囲長だけが性能を決定する指標ではない。例えば、通信路が十分に良好な場合、周囲長が6で長さ6のサイクルを1000個含む符号は、長さ4のサイクルを1個と長さ6のサイクルを2個含む符号よりも性能が悪くなる可能性がある。

SC符号の畳み込み構造は、サイクルの数がベース行列、分割行列、およびリフティング行列によってコンパクトな形で表現できることを意味しており、これにより短いサイクルの列挙と最適化が容易になる。特に、タナーグラフ内の各サイクルはプロトグラフ内の閉路からリフトされたものであり、後者はベース行列内の別の閉路から生じる。ノードの重複を許すサイクルのことを閉路と呼び、これらは「サイクル候補（cycle candidates）」とも呼ばれる。

ベース行列やプロトグラフ内のサイクル候補が、実際にプロトグラフやタナーグラフ内のサイクルを生み出すかどうかは、分割行列やリフティング行列における、そのサイクルに沿った要素の「交替和（alternating sum）」によって決定される。

\textbf{補題1（交替和条件）.}
$g \ge 2$ とし、ベース行列内の長さ $2g$ のサイクル候補 $C$ を
\[
  (j_1, i_1, j_2, i_2, \dots, j_g, i_g)
\]
とする。ここで $j_{g+1}=j_1$ であり、$(i_k, j_k)$ および $(i_k, j_{k+1})$ （$1 \le k \le g$）は分割行列 $P$ における $C$ のノードである。このとき、$C$ がプロトグラフ内で長さ $2g$ のサイクルになるための必要十分条件は以下の通りである。

NTB符号の場合：
\begin{equation}
\sum_{k=1}^g ((P)_{i_k, j_k} - (P)_{i_k, j_{k+1}}) = 0
\end{equation}

TB符号の場合：
\begin{equation}
\sum_{k=1}^g ((P)_{i_k, j_k} - (P)_{i_k, j_{k+1}}) = 0 \pmod L
\end{equation}

\subsection{SC-QLDPC符号の最適化 (Optimization of SC-QLDPC Codes)}
本節では、第4章で導入された2D-SC-HGP符号における長さ4、6、8のサイクルの最適化に焦点を当てる。2D-SC-HGP符号は、本質的にはサイズ $L_1$ と $L_2$ の2つの巡回群の直積である群上の擬アベル・リフト積符号である。

特定のサイクル候補の交替和におけるパラメータが完全に相殺される場合、$A$ と $B$ （ベース行列）が固定されれば、分割行列 $P_a$ と $P_b$ の具体的な割り当てに関わらず、サイクル条件は常に満たされる。その結果、これらのサイクルの多重度はベース行列 $A$ と $B$ のみに依存する。この性質を持つサイクルを\textbf{剛体サイクル（rigid cycles）}と呼ぶ。

同様に、その多重度が $P_a$ と $P_b$ の選択に依存するサイクルを\textbf{柔軟サイクル（flexible cycles）}と呼ぶ。我々の目的は、基となるHGPベース行列の構造とSCパラメータが与えられた下で、結果として得られるSC-HGP符号内の柔軟サイクルの総数を最小化するような $P_a$ と $P_b$ を見つけることである。

我々は、既存の確率的枠組みをSC-HGP符号に拡張することにより、分割行列を最適化する。分割行列 $P_a$ および $P_b$ の各要素が、確率分布 $p^A$ および $p^B$ に従う確率変数であると仮定し、これらを多項式と結びつけて期待値を計算する（定理4）。これにより、勾配降下法を適用して確率分布 $p^A$ と $p^B$ を結合的に最適化することが可能になる。

最適化は以下の2段階の手順で行われる。
\begin{enumerate}
    \item \textbf{GRADEアルゴリズム（勾配降下法）:} 目的関数（短いサイクルの期待値）に対する勾配を計算し、分割行列に要素を配置する確率分布 $p^A$ および $p^B$ を勾配降下法によって局所的に最適化する。
    \item \textbf{AOアルゴリズム（貪欲法）:} GRADEアルゴリズムで得られた最適な確率分布に近い経験的分布を用いて、分割行列 $P_a$ および $P_b$ を初期化する。その後、短いサイクルの重み付き数を減らすような値に要素を一つずつ置き換えていき、これ以上の改善が不可能になった時点で終了する貪欲アルゴリズム（Algorithmic Optimization）を適用する。
\end{enumerate}

この最適化手法によって、量子の可換性条件（直交性）に起因する消せないループ（剛体サイクル）の限界に近づきつつ、可能な限り短いサイクルを排除した高性能な符号の構成が可能となる。

\section{シミュレーション結果 (Simulation Results)}
6.1節では、5.2節で説明した手法を用いてSC-HGP符号を構築し、6.2節ではこれらのQLDPC符号に対するBPデコーダについて説明する。6.4節では、短いサイクルの数を最適化することの価値を示すシミュレーション結果を提示する。

\subsection{構築 (Constructions)}
4.2節で述べたように、SC-HGP符号はサイズ $r_1 \times n_1$ のベース行列 $A$、サイズ $r_2 \times n_2$ のベース行列 $B$、および結合パラメータ $m_1, m_2, L_1, L_2$ によって指定される。我々は、Algorithm 1 および Algorithm 2 を使用して、2つのSC-HGP符号のファミリを生成した。

第1のファミリには、$L_1=L_2=10$ および $(r_1, r_2, n_1, n_2) = (3, 3, 8, 8)$ を持つ7つの $[[7300, 2500]]$ 符号が含まれる。第2のファミリには、$L_1=L_2=10$ および $(r_1, r_2, n_1, n_2) = (3, 3, 7, 7)$ を持つ7つの $[[5800, 1600]]$ 符号が含まれる。各ファミリ内で、Code 2は一様分布から分割行列を抽出した結果であり、Code 4はGRADEアルゴリズムによって生成された局所的最適分布から抽出した結果である。Code 1とCode 3は、それぞれCode 2とCode 4を定義する分割行列に対してAOアルゴリズムを適用した結果である。

\subsection{BPデコーダ (BP Decoders)}
確率伝搬（BP）デコーダは、古典的なLDPC符号の復号に広く使用されている確率的デコーダである。代数的デコーダでは、エラーの数が符号の最小距離以上になると復号に失敗する。しかし、異なるエラーパターンが同一のシンドロームを導き、かつそれらの確率が異なる場合、より良い結果を得ることが可能である。BPデコーダは、変数ノード（VN）と検査ノード（CN）間のメッセージの局所的な更新の反復を通じて、各シンボルの周辺分布を効率的に推定する、最尤（ML）デコーダの低計算量な近似である。

QLDPC符号に使用されるBPデコーダでは、古典的なLDPCデコーダにおけるバイナリシンボルの尤度がパウリ行列の尤度に置き換えられる。

\subsection{縮退エラー (Degenerate Errors)}
量子エラー訂正（QEC）符号と古典的なエラー訂正符号（ECC）の大きな違いの1つは以下の通りである。古典ECCでは、いかなる $\tilde{x} \neq x$ も復号エラーにつながる。一方、量子ECCでは、$x$ と推定値 $\tilde{x}$ がQEC符号のスタビライザー群の同一のコセット（剰余類）に属している限り、$\tilde{x}$ は依然として正しい推定値となり得る。

BPデコーダが失敗したとき、推定されたベクトル $\tilde{x}$ は通常、元の $x$ に近いニア・コードワード（near-codeword）である。QEC符号には縮退エラーが存在するため、$\tilde{x}$ を近くのコードワード $x'$ に向かって押し進めることがしばしば有利となる。

本稿では、BP-OSDデコーダを採用している。しかし、長さが5000を超えるQLDPC符号に対してBP-OSDデコーダを直接実行することは、非線形な計算量のために計算コストが法外になり、エラーフロア領域からデータポイントを収集するには遅すぎる。これに対処するため、BP-OSDで直接シミュレーションを行うのではなく、以下の2段階のアプローチを採用する：
\begin{enumerate}
    \item まず、追加のBP反復を必要としない元のBP-SSF（Small-Set-Flip）デコーダを実行する。BP-SSFデコーダの復号失敗をもたらしたエラーを十分に収集し、初期のフレームエラー率（FER）を $p_e^{SSF}$ として取得する。
    \item 次に、これらのエラーのみをBP-OSDに渡し、解決されなかった復号失敗の数をさらに評価する。この未解決エラーの割合を $\alpha_e^{OSD}$ とすると、BP-OSDにおける真の論理エラー率は $p_e^{SSF}\alpha_e^{OSD}$ で与えられる。
\end{enumerate}

\subsection{数値結果 (Numerical Results)}
物理量子ビットが正しい確率が $1-p$ であり、X、Y、またはZエラーの確率が等しく $p/3$ である脱分極通信路（depolarization channel）を考慮する。我々はハイパフォーマンス・コンピューティング・クラスター上で並列にシミュレーションを行った。

最適化された符号のシミュレーション結果（図9および図10）から、以下の観察が得られた：
\begin{itemize}
    \item \textbf{短いサイクルの数を減らすことは、SC-HGP符号の性能に大きな向上をもたらす：} Code 1のCode 2に対する優位性、およびCode 3のCode 4に対する優位性は、短いサイクルの数の削減に起因している。
    \item \textbf{限られたメモリで高い性能が可能である：} メモリを $(1,1)/(1,2)$ から $(2,2)/(3,3)$ に増やすとFERは1桁改善されるが、$(2,2)$ から $(3,3)$ に増やしたときの向上はわずかである。メモリが増えると、剛体サイクル（rigid cycles）の数が柔軟サイクル（flexible cycles）の数を支配するため、柔軟サイクルの数を減らしても性能への影響は限界に達する。
    \item \textbf{剛体サイクルはGRADEとAOが性能を向上させ得る限界を決定する：} サイクル数が剛体サイクルの限界に近い場合、AOアルゴリズムによるさらなる最適化の影響は小さくなる。
    \item \textbf{性能は最小距離だけの単純な関数ではない：} 最小距離が小さいSC-HGP符号が、必ずしも最小距離が大きい符号よりも性能が劣るわけではない。復号エラー率は、最小距離だけでなく、最小重みを持つコードワードの多重度や吸収セットの存在にも依存する。
\end{itemize}

これらの結果は、我々が構築した高いレート（約1/3）を持つSC-HGP符号が、既存の低いレート（1/10未満）のHGP符号と比較して、優れた閾値と低いエラー率を達成できることを示している。

\section{結論および今後の課題}
我々は、トーリック符号を古典的な2次元空間結合（2D-SC）符号の量子的対応物として記述し、その一般化として空間結合（SC）QLDPC符号を導入した。

我々は、2次元畳み込みの枠組みを用いて、2D-SC符号のパリティ検査行列を2つの不定元を持つ多項式として表現した。我々の表現は、2D-SC符号がスタビライザー符号となるための必要十分条件である単純な代数的条件を導く。また、それは、いずれかの要素符号における短いサイクルから生じる、異なる長さの短いサイクルの数に関連する公式も導く。我々は、古典的な2D-SC符号を最適化するために用いられる手法を拡張することにより、これらの公式を用いてSCハイパーグラフ積（HGP）符号を最適化する方法を説明した。

我々のシミュレーション結果は、短いサイクルの数を減らすことがSC-HGP符号の性能に大きな向上をもたらすことを示している。我々は、限られたメモリで高い性能が可能であることを示した。

今後の課題には、ウィンドウデコーダの実装、剛体サイクル（rigid cycles）から生じるボトルネックに効率的に対処するためのデコーダの改良、およびより一般的なベース行列を用いる新しいSC-QLDPC符号の開発が含まれる。
\end{document}